\newgeometry{left=15mm,right=15mm,top=13mm,bottom=13mm,headheight=0pt,headsep=0pt,footskip=0pt}
\thispagestyle{empty}
\begingroup
\fontsize{9.3}{11.05}\selectfont
\setlength{\parskip}{3.2pt}
\setlength{\abovedisplayskip}{4pt}\setlength{\belowdisplayskip}{4pt}
\setlength{\abovedisplayshortskip}{2pt}\setlength{\belowdisplayshortskip}{2pt}
{\fontsize{18}{20}\selectfont\bfseries\color{ink}Shor's algorithm, connected}\par
{\color{accent}Order finding is phase estimation of modular multiplication. One state connects values, combs, eigenphases, and measurable frequencies.}\par

\textbf{1. Put the period in a spectrum.}
Handle even \(N\) and prime powers classically. For remaining composite \(N\), choose \(1<a<N\). A nontrivial \(\gcd(a,N)\) already gives a factor; otherwise seek \(r=\min\{d>0:a^d\equiv1\pmod N\}\). Set \(f(x)=a^x\bmod N\), \(n=\lceil\log_2N\rceil\), and \(Q=2^t\) with \(N^2\le Q<2N^2\). Register \(C\) has \(t\) qubits, \(W\) has \(n\); \(0\le x,y<Q\), \(0\le c,s<r\). The permutation \(U_a\ket w=\ket{aw\bmod N}\), extended by identity on unused labels, cycles \(\ket{e_c}=\ket{a^c\bmod N}\) with length \(r\). For \(\omega=e^{2\pi\ii/r}\),
\[
\ket{u_s}=\frac1{\sqrt r}\sum_{c=0}^{r-1}\omega^{-sc}\ket{e_c},\qquad
U_a\ket{u_s}=\omega^s\ket{u_s},\qquad
\ket1=\frac1{\sqrt r}\sum_{s=0}^{r-1}\ket{u_s}.
\]
An index shift and roots-of-unity orthogonality prove these identities. This unknown \(r\)-point Fourier basis is analytical; the circuit uses no unknown transform.

\textbf{2. Make an eigenphase relative, efficiently.}
Control makes an unobservable global phase relative, sending
\((\alpha\ket0+\beta\ket1)\ket{u_s}\) to
\((\alpha\ket0+\beta\omega^s\ket1)\ket{u_s}\). An eigenstate target stays factored. For \(x=\sum_j2^jx_j\), the ladder
\(\Lambda=\sum_x\proj x\otimes U_a^x\) uses \(t\) controlled multipliers by \(b_j=a^{2^j}\bmod N\), computed by repeated squaring. Control phases \(2^js/r\) combine to \(\omega^{sx}\), rather than independently readable bits.

\textbf{3. Use the easy input to unite three representations.}
From \(\ket0_C\ket1_W\), \(H^{\otimes t}\) gives \(Q^{-1/2}\sum_x\ket x\ket1\); then apply \(\Lambda\). Define \(M_c=\#\{x:x\equiv c\pmod r\}=1+\lfloor(Q-1-c)/r\rfloor\) and \(\ket{C_c}=M_c^{-1/2}\sum_{m=0}^{M_c-1}\ket{c+mr}\):
\[
\begin{aligned}
\ket{\Psi_2}
&=\underbrace{\frac1{\sqrt Q}\sum_x\ket x\ket{e_{x\bmod r}}}_{\text{input/output values}}
=\underbrace{\sum_c\sqrt{M_c/Q}\ket{C_c}\ket{e_c}}_{\text{orthonormal Schmidt combs}}\\[-1pt]
&=\underbrace{\frac1{\sqrt r}\sum_s\ket{\Phi_s}\ket{u_s}}_{\text{eigenphase ramps with orthogonal work tags}},
\qquad \ket{\Phi_s}=\frac1{\sqrt Q}\sum_x e^{2\pi\ii sx/r}\ket x .
\end{aligned}
\]
Values and eigenphases describe one permutation, rather than successive gates. Easy \(\ket1\) avoids unknown eigenvector preparation. For \(r>1\), the joint state is entangled.

\textbf{4. Entanglement determines the local coherences and information.}
Ignoring \(W\) gives
\[
\rho_C=\tr_W\proj{\Psi_2}
=\sum_c\frac{M_c}{Q}\proj{C_c}=\frac1r\sum_s\proj{\Phi_s},
\qquad (\rho_C)_{xx'}=\frac1Q\mathbf1_{x\equiv x'\; (\mathrm{mod}\,r)}.
\]
The diagonal is flat: measuring now gives uniform \(x\). Identical work records preserve interference within a residue class; orthogonal records remove local cross terms between classes. The Schmidt rank is \(r\), the eigenvalues are \(\lambda_c=M_c/Q\), and
\(E=S(C)=S(W)=-\sum_c\lambda_c\log_2\lambda_c\le\log_2r\) ebits, \(S(CW)=0\), \(I(C:W)=S(C)+S(W)-S(CW)=2E\), \(S(C|W)=-E\).
If \(r\mid Q\), all \(M_c=Q/r\) and \(E=\log_2r\): \(k\) Bell pairs have \(2^k\) equal branches. When \(r\nmid Q\), the ramp ensemble is not spectral: its weights \(1/r\) are not Schmidt eigenvalues. Entropy measures information in the weights, not branch count or accessible classical bits.

\textbf{5. Phase estimation reads the frequency by interference.}
Use the known Fourier matrix \((F_Q)_{xy}=Q^{-1/2}e^{2\pi\ii xy/Q}\). Apply \(F_Q^\dagger\): each \(y\) tests frequency \(y/Q\), unwinding matching phases so amplitudes add. The QPE unitary is \((F_Q^\dagger\otimes\id)\Lambda(H^{\otimes t}\otimes\id)\), and
\[
\begin{aligned}
\ket{\Psi_3}&=\frac1{\sqrt r}\sum_{y,s}D_Q(s/r-y/Q)\ket y\ket{u_s},
&D_Q(\delta)&=\frac1Q\sum_{x=0}^{Q-1}e^{2\pi\ii x\delta},\\[-1pt]
P(y)&=\bra yF_Q^\dagger\rho_CF_Q\ket y
=\frac1r\sum_s|D_Q(s/r-y/Q)|^2,
&|D_Q(\delta)|^2&=\frac{\sin^2(\pi Q\delta)}{Q^2\sin^2(\pi\delta)}.
\end{aligned}
\]
Use continuous limits at integer \(\delta\). Orthogonal \(u_s\) tags make counting statistics an exact uniform mixture of phase-estimation laws, while \(x\)-alternatives interfere within each law. This is a statistical equivalence, not hidden branch selection. If \(r\mid Q\), outputs are \(y=sQ/r\), each with probability \(1/r\); otherwise peaks broaden. The inverse QFT need not diagonalize \(\rho_C\) and preserves \(E\). Conditioned on \(y\), the work state is \(\ket{w_y}\propto\sum_sD_Q(s/r-y/Q)\ket{u_s}\), generally a superposition.

\textbf{6. Reconstruct, verify, and factor.}
For a nearest-grid output, \(|y/Q-s/r|\le1/(2Q)<1/(2r^2)\), so the reduced fraction appears among continued-fraction convergents. Such an output has probability at least \(4/\pi^2\) per phase branch. Its denominator is \(r/\gcd(s,r)\); coprime \(s\) recovers \(r\). Repeat and verify candidates: \(a^d\equiv1\pmod N\) generally proves \(r\mid d\), not minimality alone. For even \(r\) and \(z=a^{r/2}\not\equiv\pm1\pmod N\), \(N\mid(z-1)(z+1)\) and \(\gcd(z\pm1,N)\) give factors. Retry samples or bases when necessary.

\textbf{7. Preserve the right records.}
Extra scratch multiplies \((\rho_C)_{xx'}\) by \(\langle g_{x'}|g_x\rangle\). Uncompute records distinguishing inputs inside one residue class, retaining \(f(x)\). Reversing \(\Lambda\) restores the uniform input and forces \(y=0\). Measuring \(W\)'s modular value and forgetting it removes \(C|W\) entanglement but preserves \(\rho_C\), comb coherence, and Fourier statistics. Arithmetic and reconstruction cost polynomially in \(\log N\); the QFT uses \(O(t^2)\) gates.
\par\vfill
{\fontsize{8.1}{9.3}\selectfont\color{accent}Detailed derivations: explainer Sections 3--8 and 11--15. Presentation connections and its four-branch example are integrated into the full text. All identities above allow \(r\nmid Q\), except the explicitly stated exact-peak specialization.}
\clearpage
\endgroup
\restoregeometry
